
\documentclass{article}
\usepackage{colortbl}
\usepackage{makecell}
\usepackage{multirow}
\usepackage{supertabular}

\begin{document}

\newcounter{utterance}

\centering \large Interaction Transcript for game `eqbench', experiment `eq\_bench', episode 147 with reflexion{-}Qwen3.5{-}27B.
\vspace{24pt}

{ \footnotesize  \setcounter{utterance}{1}
\setlength{\tabcolsep}{0pt}
\begin{supertabular}{c@{$\;$}|p{.15\linewidth}@{}p{.15\linewidth}p{.15\linewidth}p{.15\linewidth}p{.15\linewidth}p{.15\linewidth}}
    \# & \multicolumn{2}{c}{Player} && \multicolumn{2}{c}{Game Master} \\
    \hline

    \theutterance \stepcounter{utterance}  
    & & & \multicolumn{4}{p{0.6\linewidth}}{
        \cellcolor[rgb]{0.9,0.9,0.9}{
            \makecell[{{p{\linewidth}}}]{
                \texttt{\tiny{[P1$\langle$GM]}}
                \\ 
\texttt{Your task is to predict the likely emotional responses of a character in this dialogue:} \\
\\ 
\texttt{Jackson: Sanchez. Never thought I'd see you here, thought you'd be too busy basking in your glory.} \\
\texttt{Alejandro: Riley, always the charmer. I could say the same about you. How's the knee?} \\
\texttt{Jackson: Better than your conscience, I reckon.} \\
\texttt{Alejandro: Ah, the ever{-}righteous Riley. You think you're the only one who's had to fight for everything?} \\
\texttt{Jackson: I don't recall stepping on others to get to the top.} \\
\texttt{Alejandro: And I don't recall asking for your opinion.} \\
\texttt{Jackson: You never did. Too busy making a spectacle of yourself.} \\
\texttt{Alejandro: And you're too busy playing the victim. } \\
\texttt{Jackson: I ain't the one parading around with another man's woman.} \\
\texttt{Alejandro: And I ain't the one who couldn't keep her.} \\
\texttt{Jackson: You think this is about her? This is about respect, something you wouldn't understand.} \\
\texttt{Alejandro: And you think this is about respect? This is about survival, something you've never had to worry about.} \\
\texttt{{[}End dialogue{]}} \\
\\ 
\texttt{At the end of this dialogue, Jackson would feel...} \\
\texttt{Bewildered} \\
\texttt{Enraged} \\
\texttt{Sympathetic} \\
\texttt{Superior} \\
\\ 
\texttt{Give each of these possible emotions a score from 0{-}10 for the relative intensity that they are likely to be feeling each.} \\
\\ 
\texttt{You must output in the following format, including headings (of course, you should give your own scores), with no additional commentary:} \\
\\ 
\texttt{Bewildered: <score>} \\
\texttt{Enraged: <score>} \\
\texttt{Sympathetic: <score>} \\
\texttt{Superior: <score>} \\
\\ 
\\ 
\texttt{{[}End of answer{]}} \\
\\ 
\texttt{Remember: zero is a valid score, meaning they are likely not feeling that emotion. You must score at least one emotion > 0.} \\
\\ 
\texttt{Your answer:} \\
            }
        }
    }
     \\ \\

    \theutterance \stepcounter{utterance}  
    & \multicolumn{4}{p{0.6\linewidth}}{
        \cellcolor[rgb]{0.561,0.737,0.561}{
            \makecell[{{p{\linewidth}}}]{
                \texttt{\tiny{[P1$\rangle$GM]}}
                \texttt{Bewildered: 1} \\
\texttt{Enraged: 9} \\
\texttt{Sympathetic: 0} \\
\texttt{Superior: 4} \\
            }
        }
    }
    & & \\ \\

    \theutterance \stepcounter{utterance}  
    & & & \multicolumn{2}{p{0.3\linewidth}}{
        \cellcolor[rgb]{0.9,0.9,0.9}{
            \makecell[{{p{\linewidth}}}]{
                \texttt{\tiny{[GM$|$GM]}}
                \texttt{Bewildered: 1} \\
\texttt{Enraged: 9} \\
\texttt{Sympathetic: 0} \\
\texttt{Superior: 4} \\
            }
        }
    }
    & & \\ \\

    \theutterance \stepcounter{utterance}  
    & & & \multicolumn{2}{p{0.3\linewidth}}{
        \cellcolor[rgb]{0.9,0.9,0.9}{
            \makecell[{{p{\linewidth}}}]{
                \texttt{\tiny{[GM$|$GM]}}
                \texttt{Bewildered: 0} \\
\texttt{Enraged: 8} \\
\texttt{Sympathetic: 0} \\
\texttt{Superior: 4} \\
            }
        }
    }
    & & \\ \\

    \theutterance \stepcounter{utterance}  
    & & & \multicolumn{2}{p{0.3\linewidth}}{
        \cellcolor[rgb]{0.9,0.9,0.9}{
            \makecell[{{p{\linewidth}}}]{
                \texttt{\tiny{[GM$|$GM]}}
                \texttt{game\_result = LOSE} \\
            }
        }
    }
    & & \\ \\

    \theutterance \stepcounter{utterance}  
    & & \multicolumn{4}{p{0.6\linewidth}}{
        \cellcolor[rgb]{0.95,0.95,0.95}{
            \makecell[{{p{\linewidth}}}]{
                \texttt{\tiny{[P1$\leftrightarrow$P]}}
                \texttt{Post{-}episode Reflexion (not sent to GM):} \\
\texttt{New plan: My previous attempt failed because I assigned a non{-}zero score to "Superior" (4), which contradicts the dialogue's context where Jackson is being aggressively attacked and defensive, not feeling superior. The dialogue is a heated exchange of insults where Jackson is clearly provoked ("Never thought I'd see you here," "Better than your conscience," "parading around with another man's woman"). The most dominant emotion is "Enraged," which I correctly identified as high, but I should have recognized that "Superior" is inconsistent with Jackson's reactive and defensive stance. In the future, I will strictly evaluate each emotion against the character's specific lines: if a character is being insulted and retaliating defensively, "Superior" should be 0. I will also ensure "Bewildered" remains low (as the conflict is clear) and "Sympathetic" is 0 (as there is no empathy shown). The corrected strategy is to assign high intensity only to "Enraged" and 0 to all others, or very low scores to secondary emotions that are truly plausible, avoiding any score for emotions that contradict the character's defensive posture. Specifically, I will set Superior to 0.} \\
            }
        }
    }
    & \\ \\

\end{supertabular}
}

\end{document}
